\scanpage{103}
\mspage{\textcircled{6}}
5.6. Partons de $R\in\mathrm{Ob}\,\mathrm{Cat}_{\lambda}$, avec $I\xrightarrow{\varphi\,\unsure{\subset}} \mathrm{Ob}\,R$ tels
que $R$ soit $\lambda$-engendrée par $I$. Considérons la cat.
2-\deleted{\unsure{co}}\inserted{\unsure{co}}fibrée $T$ sur $(\mathrm{Cat}_{\lambda})$ définie par $R$, et sa structure
supplémentaire $T\xrightarrow{\varphi}\tau^I$ déduite de $\varphi$. \unsure{Prouvons|Prouverons|Prouver}
que $(T,\varphi)$ est un $\lambda$-$I$-type ? (Il faut peut-être
\margin{$|\!|\!|$}renforcer la notion de sous-catégorie $\lambda$-engendrée, en
tenant compte des flèches, pas seulement des objets... ?)

5.7. Si on laisse tomber $I$, on \deleted{trouve que}\inserted{doit présumer}
Catégorie 2-cofibrée $T$ sur $(\mathrm{Cat}_{\lambda})$ est un $\lambda$-type, si c'est
un $\lambda$-$I$-type pour \inserted{$I$ et} $\varphi:T\longrightarrow\tau^I$ convenables. Cela
signifie donc (\unsure{en|si} ce qui précède \unsure{généraliser|généralise}) que
$T$ est 2-représentable par un $R\in\mathrm{Ob}(\mathrm{Cat}_{\lambda})$, où
$\mathrm{Ob}\,R$ admet une \deleted{petite} partie qui soit $\lambda$-génératrice
(en un sens convenable à dégager).

% slanted
\margin{\uline{Question} Si $\underset{\leftarrow}{\lambda}$ et $\underset{\rightarrow}{\lambda}$ ne comportent que des diagrammes finis, en nb fini, alors une $R\in\mathrm{Ob}\,\mathrm{Cat}_\lambda$ muni d'un $I\longrightarrow\mathrm{Ob}\,R$ \unsure{provenant} d'un $\lambda$-type, est-elle une \uline{petite} catégorie ?}

\fbox{6} \uline{$\lambda$-variable} on devrait trouver \inserted{pour $\lambda\subset\lambda'$} un 2-foncteur
\[
(\lambda\text{-types})\xrightarrow{\unsure{\gamma|\eta}_{\lambda,\lambda'}}(\lambda'\text{-types})\tag*{(6.1)}
\]
\margin{$(\gamma_{\lambda,\lambda'}(T_\lambda))(C)=T_\lambda(C)$ si $C\in\mathrm{Ob}\,\mathrm{Cat}_{\lambda'}$ $(\underset{\mathrm{Ob}\,\mathrm{Cat}_{\lambda}}{\cap})$}
qui corresponde à un 2-foncteur
\[
(\mathrm{Cat}_{\lambda})\xrightarrow{\overline\gamma_{\lambda,\lambda'}}(\mathrm{Cat}_{\lambda'})\tag*{(6.2)}
\]
``$\lambda'$-catégorie enveloppante''. Il faudrait donner des
critères moyennant lesquels (6.1) est pl. fidèle.
\[
\begin{aligned}
[\ \underline{\text{Ex}}\quad\lambda&=\varprojlim\text{ finies},\quad\lambda'\ \unsure{\subset|=}\ \varprojlim\text{ quelc.}\quad(\overline\gamma(C)=\mathrm{Pro}(C)\,?)\\
\lambda&=\varinjlim\text{ finies},\quad\lambda'=\varinjlim\text{ quelc.}\quad(\overline\gamma(C)=\mathrm{Ind}(C)\,?)]
\end{aligned}
\]
\margin{? $||$} $\gamma_{\lambda,\lambda'}$ est-il tjrs \deleted{\unsure{pl}}-fidèle ?, i.e
$\underline{\mathrm{Hom}}_\lambda(T_1,T_2)\longrightarrow\mathrm{Hom}_{\lambda'}(T_1,T_2)$
est-il tjrs une équivalence pour $T_1$ et $T_2$ deux $\lambda$-types ? Sans
doute faux, mais ce doit être pl. fidèle si $\underset{\rightarrow}{\lambda}=\varnothing$,

\fbox{7} \uline{Dualité}. \uline{Opposé} d'un ``type'' $\lambda=(\underset{\leftarrow}{\lambda},\underset{\rightarrow}{\lambda})$ : on échange
les types de diagrammes utilisés pour des $\varprojlim$ resp $\varinjlim$.
Alors pour tt $\lambda$-type $T_{\deleted{\unsure{\lambda}}}$, on définit un $\lambda^o$-type $T^o$ \unsure{t.q}
\[
T^o_{\deleted{\unsure{\lambda}}}(C)\simeq T_{\deleted{\unsure{\lambda}}}(C^o)^o\tag*{(\unsure{7.1|6.3})}
\]
% archive sheet number: 102
